% !PRINTED 486
\weberpage{486}
%<<CONT
\[
\begin{aligned}
y &= \psi(\alpha, \alpha_1 \ldots \alpha_{r-1}) \\
y_1 &= \psi(\beta, \beta_1 \ldots \beta_{r-1}) \\
& \phantom{=} \; \cdots\cdots\cdots\cdots \\
y_{s-1} &= \psi(\sigma, \sigma_1 \ldots \sigma_{r-1})
\end{aligned}
\tag{8}
\]
alle von einander verschieden sind.

Um die Möglichkeit hiervon einzusehen, können wir z.\,B.
\[
\begin{aligned}
\psi(t, A) &= (t - \alpha)(t - \alpha_1) \ldots (t - \alpha_{r-1}), \\
\psi(t, B) &= (t - \beta)(t - \beta_1) \ldots (t - \beta_{r-1}), \\
& \phantom{=} \; \cdots\cdots\cdots\cdots \\
\psi(t, S) &= (t - \sigma)(t - \sigma_1) \ldots (t - \sigma_{r-1})
\end{aligned}
\tag{9}
\]
setzen, und dann lässt sich (nach \S.\,143, 1.) für $t$ ein solcher
rationaler Werth finden, dass diese Grössen alle von einander
verschieden ausfallen. Diese Werthe können also für die $y$ in (8)
genommen werden.

Wenn wir nun unter $u$ eine unabhängige Variable verstehen und
\weq{10}{\varphi(u) = (u - y)(u - y_1) \ldots (u - y_{s-1})}
setzen, so ist $\varphi(u)$, da seine Co\"efficienten durch die Permutationen
der Gruppe ungeändert bleiben, eine Function in $\Omega$, deren Wurzeln die Grössen $y$ sind. \gesperrt{Sie ist überdies irreducibel}, denn
aus der vorausgesetzten Transitivität der Gruppe folgt, dass,
wenn eine rationale Function von $y$ für einen der Werthe (8)
verschwindet, sie auch für alle anderen verschwinden muss.

Ist $\omega$ irgend eine symmetrische Function der $\alpha, \alpha_1 \ldots \alpha_{r-1}$,
so schliessen wir aus der Betrachtung des Ausdruckes
\[
\varphi(u) \left( \frac{\omega}{u - y} + \frac{\omega_1}{u - y_1} + \cdots + \frac{\omega_{s-1}}{u - y_{s-1}} \right),
\]
der eine ganze rationale Function von $u$ in $\Omega$ ist, ganz wie oben,
indem wir $u = y$ setzen, dass $\omega$ rational durch $y$ dargestellt
werden kann, und demnach kann auch die Function
\weq{11}{\psi(t, A) = \psi(t, y),}
deren Wurzeln die Grössen $\alpha, \alpha_1 \ldots \alpha_{r-1}$ sind, rational durch
$y$ ausgedrückt werden.

\gesperrt{Auch die Gleichung $\psi(t, A) = 0$ ist irreducibel
in dem Körper $\Omega(y)$}; denn ist $\psi_1(t, y)$ ein rationaler Theiler
von $\psi(t, y)$, und ist
%CONT>>
